Wir können von der lokalen Situation ausgehen, dass auf dem trivialen Vektorbündel
\maabbdisp {} {U \times \R} {U
} {}
ein Zusammenhang vorliegt, wobei wir

\mavergleichskette
{\vergleichskette
{ U
}
{ \subseteq }{  \R^n
}
{  }{
}
{    }{
}
{   }{
}
}
{}{}{}
als
\definitionsverweis {sternförmig}{}{}
annehmen können. Der Zusammenhang wird durch die
\definitionsverweis {Christoffelsymbole}{}{}
$\Gamma_1 , \ldots , \Gamma_n$ beschrieben. Da der Zusammenhang als krümmungsfrei vorausgesetzt wird, gilt nach
[[Reelles Vektorbündel/Rang 1/Offene Menge/Krümmungsoperator/Fakt|Fakt]],
dass die
$1$-\definitionsverweis {Differentialform}{}{}
$\sum_{i = 1}^n\Gamma_i dx_i$
\definitionsverweis {geschlossen}{}{}
ist. Nach
[[1-Form/Vektorfeld/Geschlossen gdw Integrabilitätsbedingung/Aufgabe|Aufgabe]]
bedeutet dies, dass das zugehörige Vektorfeld
\maabbdisp {\left( \Gamma_1   , \,   \ldots  , \,  \Gamma_n                 \right)} { U } {\R^n
} {}
die
\definitionsverweis {Integrabilitätsbedingung}{}{}
erfüllt. Dies bedeutet nach
[[Teilmenge/R^n/Sternförmig/Gradientenfeld/Charakterisierung/Fakt|Fakt]],
dass auf der sternförmigen offenen Menge $U$ ein
\definitionsverweis {Gradientenfeld}{}{}
vorliegt. Es gibt also eine stetig differenzierbare Funktion
\maabbdisp {f} { U } { \R
} {}
mit

\mavergleichskettedisp
{\vergleichskette
{ \Gamma_i
}
{ =} { \partial_i f
}
{ } {
}
{ } {
}
{ } {
}
}
{}{}{}
für alle $i$. Wir setzen

\mavergleichskettedisp
{\vergleichskette
{ g
}
{ =} { e^f
}
{ } {
}
{ } {
}
{ } {
}
}
{}{}{}
und fassen $g$ als einen Schnitt im trivialen Bündel $U \times \R$ auf. Dabei ist nach
[[Differenzierbare Mannigfaltigkeit/R/Vektorbündel/Linearer Zusammenhang/Lokale Beschreibung/Fakt|Fakt]]

\mavergleichskettedisp
{\vergleichskette
{ \nabla_{\partial_i} g
}
{ =} { g \Gamma_i -  \partial_i g
}
{ =} { e^f \Gamma_i - \partial_i e^f
}
{ =} { e^f \Gamma_i - e^f \partial_i f
}
{ =} { 0
}
}
{}{}{,}
d.h. dass $g$ ein nullstellenfreier
\definitionsverweis {horizontaler Schnitt}{}{}
ist, der $\nabla$ als lokal integrabel erweist.

 [[Kategorie:Latexseite]]
